\chapter{Marchés au comptant}\label{ch:m3:spot-markets}

Le même bitcoin se négocie sur des dizaines de plateformes, à des dizaines de prix. La plupart du temps, les prix diffèrent
de quelques points de base, soit la largeur des frais~; parfois, ce n'est pas le cas. De décembre 2017 à début février
2018, le bitcoin acheté en won à Séoul coûtait en moyenne plus de 15~\% de plus que le bitcoin
acheté en dollars aux États-Unis, et pendant plusieurs jours, 40~\% de plus. N'importe qui pouvait voir l'écart, les coins
pouvaient traverser le monde en une heure, et l'écart est resté ouvert pendant des semaines. Ce chapitre porte sur la façon dont
les prix au comptant des crypto-actifs se forment sur des plateformes fragmentées, sur les paires dans lesquelles ils sont cotés, sur les arbitrages qui
les relient et sur ce qui limite ces arbitrages, au premier rang desquels figurent le temps et l'autorisation nécessaires pour déplacer de l'argent, et enfin
sur la part réelle des volumes que déclarent les plateformes.

\section{Fragmentation et formation des prix entre plateformes}

Un crypto-actif n'a pas de marché primaire. Il n'existe ni bourse de cotation dont l'enchère fixe le prix de référence, ni flux consolidé, ni règle
qui achemine un ordre vers le meilleur prix, comme sur les marchés actions du livre~1, chapitre~9. Chaque \omterm{def:m3:centralised-exchanges:cex}{plateforme centralisée} (CEX) fait tourner
son propre carnet d'ordres sur ses propres paires, pour ses propres clients, tandis que les plateformes décentralisées ajoutent leurs pools
(\cref{ch:m3:automated-market-makers}), et le prix du bitcoin est celui auquel
les arbitragistes parviennent à aligner ces carnets.

Ils s'alignent étroitement là où le capital d'arbitrage circule librement. Makarov et Schoar, qui ont étudié les données tick par tick de 34 bourses dans 19
pays, ont constaté que les écarts de prix entre bourses d'un même pays ne dépassaient typiquement
pas 1~\% en moyenne, tandis que les écarts entre pays étaient bien plus grands, et que, entre crypto-actifs sur les mêmes plateformes,
ils étaient encore plus faibles~: lorsque le bitcoin en won coûtait plus de 20~\% de plus
que le bitcoin en dollars, le prix de l'ether exprimé en bitcoin différait d'environ 3~\% en moyenne entre les deux
pays. Les coins circulaient librement~; la monnaie fiduciaire, non. Le prix d'un crypto-actif se pense donc au mieux
comme un prix par zone monétaire, reliés par le coût du transfert de fonds entre elles.

Au sein d'une zone monétaire, la formation des prix se produit là où l'information s'échange~: dans les carnets les plus profonds, et dans les
contrats perpétuels (\cref{ch:m3:perpetual-futures}) qui se négocient aux côtés du marché au comptant sur les mêmes plateformes.
Les petites plateformes suivent, leurs prix étant maintenus en ligne par des traders qui cotent sur elles en se couvrant sur les grandes.
Pour une société de trading, les questions pratiques sont de savoir où le prix se forme, à quelle vitesse les autres suivent,
et ce qu'il en coûte de négocier sur chaque plateforme et de déplacer l'inventaire de l'une à l'autre.

\section{Paires cotées en stablecoins}

Les paires cryptographiques suivent les conventions du marché des changes (livre~2, chapitre~14)~: dans BTC/USDT, le bitcoin
est la devise de base et l'USDT la devise de cotation, et le prix est le nombre d'USDT par bitcoin.
La plupart des plateformes offshore cotent leurs carnets les plus profonds contre des \omterm{def:m3:blockchains-for-traders:stable}{stablecoins} plutôt que contre des dollars, car une plateforme
sans relations bancaires peut détenir des \omterm{def:m3:blockchains-for-traders:stable}{stablecoins} mais pas de dépôts en dollars, et parce que les \omterm{def:m3:blockchains-for-traders:stable}{stablecoins} circulent entre plateformes
via les chaînes à toute heure.

\begin{dated}{2026-09}{Devises de cotation sur une grande plateforme}\label{dat:m3:spot-markets:quotes}
Le 24 septembre 2026, le point d'accès aux informations sur les échanges de Binance recensait 1\,372 paires au comptant en négociation, dont
495 cotées en USDT, 264 en USDC, 309 en livres turques, 38 en bitcoin et 29 en euros~; seules quelques-unes
étaient cotées en dollars.
\end{dated}

Un prix en USDT n'est un prix en dollars que si l'USDT vaut un dollar. Un trader qui compare BTC/USDT sur une plateforme
avec BTC/USD sur une autre compare deux choses différentes, reliées par le cours USDT/USD, qui est lui-même coté
et peut bouger~: de quelques points de base en temps normal, bien davantage lors d'un \omterm{def:m3:blockchains-for-traders:stable}{décrochage} (depeg)
(\cref{ch:m3:blockchains-for-traders}). Le prix BTC/USDT et le prix BTC/USD donnent un cours croisé USDT/USD implicite, et chaque arbitrage entre les deux carnets
est une position sur ce cours croisé. Une société de trading tient ses comptes dans une seule devise
et valorise ses soldes en \omterm{def:m3:blockchains-for-traders:stable}{stablecoins} à leur cours de marché, et non au pair.

\section{Arbitrage triangulaire et arbitrage entre plateformes}

Deux arbitrages maintiennent les prix alignés.

\begin{definition}[Arbitrage triangulaire]\label{def:m3:spot-markets:tri}
Un \emph{arbitrage triangulaire}\index{arbitrage triangulaire} est un cycle de trois transactions sur une même plateforme qui part d'un actif
et y revient en passant par deux autres, par exemple USDT vers bitcoin, bitcoin vers ether puis ether vers USDT, exécuté lorsque le produit
des trois cours de change, frais déduits, dépasse un.
\end{definition}

\begin{proposition}[Bande d'absence d'arbitrage d'un cours croisé]\label{prop:m3:spot-markets:band}
Soit une plateforme qui cote $A/Q$, $B/Q$ et $B/A$, et qui prélève des frais preneurs $f$ sur chaque transaction. Le cycle
$Q \to A \to B \to Q$ est rentable si et seulement si
\[
\frac{\text{cours acheteur}_{B/Q}}{\text{cours vendeur}_{A/Q}\,\text{cours vendeur}_{B/A}}\,(1 - f)^3 > 1,
\]
et le cycle inverse $Q \to B \to A \to Q$ est rentable si et seulement si
$\text{cours acheteur}_{B/A}\,\text{cours acheteur}_{A/Q}\,(1 - f)^3 / \text{cours vendeur}_{B/Q} > 1$. Aucun des deux n'est rentable tant que le cours
$B/A$ coté reste dans une bande autour du cours croisé implicite $B/Q \div A/Q$, dont la demi-largeur vaut environ $3f$
plus les fourchettes.
\end{proposition}

\begin{proof}
Dépensons une unité de $Q$~: l'achat de $A$ à son cours vendeur (ask) donne $(1-f)/\text{ask}_{A/Q}$ unités de $A$, et l'achat de $B$ avec celles-ci
au cours vendeur de $B/A$ donne $(1-f)^2/(\text{ask}_{A/Q}\,\text{ask}_{B/A})$ unités de $B$, et la vente de $B$ à son cours acheteur restitue le produit annoncé.
Le cycle inverse relève du même calcul dans l'autre sens. En prenant les logarithmes, chaque condition dit que le logarithme du cours croisé coté
diffère de celui du cours croisé implicite d'une quantité supérieure à
$-3\ln(1-f) \approx 3f$, plus les demi-fourchettes.
\end{proof}

À des frais preneurs de 10 points de base, la bande fait environ 30 points de base de large de chaque côté, et les teneurs de marché propres à une plateforme
maintiennent ses cours croisés bien à l'intérieur de celle-ci~: sur une grande plateforme isolée, les opportunités triangulaires sont rares et
petites, et captées par qui dispose des frais les plus bas et de la connexion la plus rapide. Le tutoriel les compte dans des carnets synthétiques~:
1\,254 cycles rentables avant frais sur 500 instantanés de trois plateformes, dont 26 à des frais preneurs
de 10 points de base.

\begin{omfigure}
\centering
\begin{tikzpicture}[font=\small,
  c/.style={draw,thick,circle,minimum size=1.3cm,align=center},
  arr/.style={-{Stealth[length=2.2mm]},thick}]
\node[c,draw=omDef,fill=omDef!8] (q) at (0,0) {USDT};
\node[c,draw=omThm,fill=omThm!8] (a) at (3.6,2.2) {BTC};
\node[c,draw=omProp,fill=omProp!8] (b) at (7.2,0) {ETH};
\draw[arr] (q) -- node[above left,font=\footnotesize,align=center] {acheter BTC/USDT\\au ask} (a);
\draw[arr] (a) -- node[above right,font=\footnotesize,align=center] {acheter ETH/BTC\\au ask} (b);
\draw[arr] (b) -- node[below,font=\footnotesize] {vendre ETH/USDT au bid} (q);
\node[font=\footnotesize,align=center] at (3.6,0.75) {trois frais preneurs~:\\rentable si le produit des\\cours dépasse $(1-f)^{-3}$};
\end{tikzpicture}

\omcaption{Un cycle triangulaire sur une plateforme~: USDT vers bitcoin, bitcoin vers ether, ether vers USDT.
Le cycle, ou son inverse, n'est rentable que lorsque le cours ETH/BTC coté sort de la bande autour du cours croisé
implicite, bande fixée par trois frais preneurs et par les fourchettes. Schéma.}\label{fig:m3:spot-markets:triangle}
\end{omfigure}

\begin{definition}[Arbitrage entre plateformes, inventaire prépositionné]\label{def:m3:spot-markets:cross}
Un \emph{arbitrage entre plateformes}\index{arbitrage entre plateformes} consiste à acheter un actif sur une plateforme et à le vendre sur une autre
où il est plus cher. Un \emph{inventaire prépositionné}\index{inventaire prépositionné} est le stock de
l'actif et de la devise de cotation qu'un arbitragiste conserve sur chaque plateforme afin que les deux jambes puissent être exécutées
simultanément, l'inventaire étant rapatrié entre plateformes plus tard, lorsque cela coûte le moins cher.
\end{definition}

Avec un inventaire sur les deux plateformes, l'arbitragiste vend sur la plateforme chère à partir de son stock qui s'y trouve et achète
sur la plateforme bon marché avec ses liquidités qui s'y trouvent, au même moment~; la transaction est verrouillée aux prix observés. Les coûts
sont deux frais preneurs, le coût éventuel du rééquilibrage (frais de retrait, frais de réseau et durée de transit de l'inventaire),
ainsi que le capital immobilisé~: un inventaire déposé sur les plateformes est exposé à leur défaillance
(\cref{ch:m3:centralised-exchanges}). Sans inventaire, l'arbitragiste doit acheter, transférer, puis
vendre, et le prix peut évoluer pendant que les \omterm{def:m3:blockchains-for-traders:key}{jetons} sont en transit.

\section{Latence de transfert et risque}

\begin{definition}[Latence de transfert]\label{def:m3:spot-markets:latency}
Une \emph{latence de transfert}\index{latence de transfert} est le délai entre la décision de déplacer un actif d'une
plateforme à une autre et le moment où il peut être négocié sur la plateforme de destination~: le traitement du retrait
par la plateforme d'origine, l'inclusion dans la chaîne et les confirmations exigées par la destination, ainsi que
l'inscription du crédit sur le compte de destination.
\end{definition}

La part de la chaîne dépend de ses temps de bloc et de sa \omterm{def:m3:blockchains-for-traders:finality}{finalité} (\cref{ch:m3:blockchains-for-traders})~; celle des
plateformes, souvent plus longue, dépend de leurs propres processus, de leurs contrôles et des soldes de leurs portefeuilles chauds. Pour un arbitragiste
qui ne dispose pas d'inventaire à destination, la \omterm{def:m3:spot-markets:latency}{latence de transfert} transforme un écart verrouillé en pari.

\begin{proposition}[Probabilité qu'un arbitrage par transfert soit perdant]\label{prop:m3:spot-markets:ploss}
Acheter à un prix $P$ lorsque le prix sur l'autre plateforme est supérieur d'un écart de $g$ (en fraction de $P$), avec des
coûts $c$, puis vendre après un transfert de durée $t$ (en années), pendant lequel le logarithme du prix varie d'une
quantité normale de volatilité $\sigma$ et sans dérive. Au premier ordre, la transaction est perdante avec la probabilité
\[
\Pr(\text{perte}) \approx \Phi\!\left(-\frac{g - c}{\sigma\sqrt{t}}\right),
\]
où $\Phi$ est la fonction de répartition de la loi normale centrée réduite.
\end{proposition}

\begin{proof}
Le rendement net vaut $g - c + \varepsilon$ avec $\varepsilon \sim \mathcal{N}(0, \sigma^2 t)$ au premier ordre
en les petites quantités~; il est négatif lorsque $\varepsilon < -(g - c)$.
\end{proof}

\begin{omfigure}
\begin{tikzpicture}
\begin{axis}[width=11cm,height=5cm,
  xlabel={temps de transfert, minutes},ylabel={probabilité de perte, \%},
  tick label style={font=\small},label style={font=\small},
  xmin=0,xmax=60,ymin=0,ymax=35,grid=major,grid style={gray!25},
  legend columns=3,legend style={font=\footnotesize,at={(0.5,-0.3)},anchor=north,draw=none,fill=none,/tikz/every even column/.append style={column sep=8pt}},
  table/col sep=comma]
\addplot[omThm,very thick] table[x=minutes,y=g50]{figdata/markets-3/16-spot-markets/latency.csv};\addlegendentry{écart de 50 pb}
\addplot[omDef,very thick,dashed] table[x=minutes,y=g100]{figdata/markets-3/16-spot-markets/latency.csv};\addlegendentry{écart de 100 pb}
\addplot[omExo,very thick,dotted] table[x=minutes,y=g200]{figdata/markets-3/16-spot-markets/latency.csv};\addlegendentry{écart de 200 pb}
\end{axis}
\end{tikzpicture}

\omcaption{Probabilité qu'un \omterm{def:m3:spot-markets:cross}{arbitrage entre plateformes} exécuté en transférant l'actif perde de l'argent, en fonction
du temps de transfert, pour trois écarts, avec 20 points de base de coûts et une volatilité annuelle de 60~\%
(\cref{prop:m3:spot-markets:ploss}). Un écart de 50 points de base perd environ une fois sur trois au bout d'une heure.
Paramètres illustratifs. Données~: le tutoriel du chapitre.}\label{fig:m3:spot-markets:latency}
\end{omfigure}

\Cref{fig:m3:spot-markets:latency} explique pourquoi des écarts de quelques dizaines de points de base subsistent entre
des plateformes lentes à créditer les dépôts~: avec une volatilité annuelle de 60~\%, l'écart type du prix sur une heure
vaut environ 64 points de base, soit davantage que l'écart. La réponse professionnelle consiste à disposer d'un \omterm{def:m3:spot-markets:cross}{inventaire
prépositionné} sur chaque plateforme qui compte, à appliquer un processus de rééquilibrage qui le déplace lorsque les frais et la
congestion sont faibles, et à couvrir l'inventaire lui-même sur une plateforme de contrats à terme. La construction de ce
chapitre recherche ces deux types d'opportunités, en imputant à une transaction par transfert le risque de prix lié à sa latence, et à une transaction issue de l'inventaire le coût du rééquilibrage.

\section{La prime kimchi}

\begin{definition}[Prime kimchi]\label{def:m3:spot-markets:kimchi}
La \emph{prime kimchi}\index{prime kimchi} est l'excédent du prix en won d'un crypto-actif sur les bourses
sud-coréennes par rapport à son prix ailleurs, converti en won au taux de change du marché.
\end{definition}

Makarov et Schoar ont constaté que, pendant de larges parties de 2017 et du début de 2018, le bitcoin sur les bourses coréennes
se négociait plus de 20~\% au-dessus du prix aux États-Unis, le plus élevé de tous~; que le rapport de prix moyen journalier entre
la Corée et les États-Unis, de décembre 2017 au début de février 2018, dépassait 15~\% et atteignait 40~\% pendant plusieurs jours~;
et que les primes s'ouvraient lors des fortes hausses du bitcoin. Ils ont estimé les profits d'arbitrage disponibles
entre la Corée et les États-Unis sur quatre mois à 1.275~milliards de dollars.

La prime a persisté parce que l'arbitrage comporte une jambe en monnaie fiduciaire. Les \omterm{def:m3:blockchains-for-traders:key}{jetons} peuvent être achetés à l'étranger, envoyés en Corée
et vendus contre des won en une heure~; les won doivent ensuite être convertis en dollars et renvoyés à l'étranger pour boucler
l'aller-retour et recommencer. Cette jambe était limitée~: les bourses coréennes servaient des résidents disposant d'un compte
bancaire local, et les résidents qui transféraient de fortes sommes à l'étranger devaient les justifier. Les mêmes auteurs notent
que les bourses coréennes et japonaises n'autorisaient pas la vente à découvert, ce qui écartait l'autre moyen d'arbitrer la prime, vendre d'abord et livrer ensuite.

\begin{dated}{2026-09}{Justificatifs de transferts sortants en Corée}\label{dat:m3:spot-markets:korea}
Makarov et Schoar (2020) décrivent qu'en Corée les résidents et les sociétés qui font sortir plus de 50\,000~USD
du pays en une année doivent soumettre aux autorités des justificatifs prouvant les motifs des transferts,
lesquels ne sont pas toujours approuvés, et que les bourses opérant dans plusieurs pays n'ouvrent des comptes aux citoyens d'un pays que
dans leur monnaie locale. À partir du 30 janvier 2018, la Commission des services financiers de Corée n'autorisa le trading de crypto-actifs que
par des comptes bancaires au nom réel, ouverts auprès de la banque de la bourse, ferma les nouveaux comptes aux étrangers et
aux mineurs, et obligea les banques à traiter comme suspects les dépôts ou retraits dépassant 10~millions de KRW par jour. Les règles et les seuils évoluent~;
une société vérifie les règles en vigueur auprès de ses conseils avant de trader.
\end{dated}

Un contrôle des capitaux n'interdit pas l'arbitrage~; il le tarife. Le problème du week-end calcule la prime en deçà de
laquelle l'aller-retour perd de l'argent, compte tenu des frais, du risque de prix du transfert et du coût de sortie des won.

\section{Wash trading et mesure de la liquidité réelle}

Le volume sert à classer les plateformes, et le rang d'une plateforme attire les admissions à la cote et les clients. Beaucoup de plateformes ont publié un
volume qui n'a pas eu lieu.

\begin{definition}[Wash trading]\label{def:m3:spot-markets:wash}
Le \emph{wash trading}\index{wash trading} est une négociation dans laquelle le même bénéficiaire effectif, ou des parties agissant de
concert, se trouvent des deux côtés, de sorte qu'aucun risque ne change de mains~; il gonfle le volume publié et peut induire
en erreur sur l'activité et le prix d'un marché.
\end{definition}

En mars 2019, la société de gestion Bitwise a présenté au personnel de la SEC une analyse soutenant qu'environ 95~\% des
quelque 6~milliards de dollars de volume quotidien de bitcoin publiés par les agrégateurs de données relevaient
d'un \omterm{def:m3:spot-markets:wash}{wash trading} factice ou sans substance économique, et que le marché réel se concentrait sur dix bourses dont les prix évoluaient
de concert. Cong, Li, Tang et Yang ont ensuite appliqué trois tests statistiques à 29 bourses. Les transactions sur les bourses
régulées présentaient les schémas observés partout~: premiers chiffres significatifs des tailles de transactions suivant la loi de Benford, regroupement des
tailles sur des nombres ronds et distributions des tailles à queues épaisses. Les bourses non régulées échouaient souvent à ces tests, et les
auteurs ont estimé que le \omterm{def:m3:spot-markets:wash}{wash trading} représentait en moyenne plus de 70~\% du volume publié des bourses non régulées.

Les deux tests implémentés par le tutoriel sont simples. La loi de Benford énonce que, dans de nombreuses données d'origine naturelle, le
premier chiffre significatif vaut $d$ avec la probabilité $\log_{10}(1 + 1/d)$~: un 1 environ 30~\% du temps, un 9 moins de 5~\%.
Les tailles produites par une machine et tirées uniformément dans un intervalle ne la suivent pas. Et les acteurs humains échangent des tailles rondes, 0.01 ou 0.1
bitcoin, bien plus souvent que 0.0137~; un robot qui émet du volume ne s'en soucie pas. Cong et ses coauteurs ont pris comme référence le rapport entre transactions non arrondies et transactions arrondies sur les
bourses régulées~: les transactions arrondies d'une autre bourse, multipliées par ce
rapport, estiment ses transactions authentiques, et l'excédent de transactions non arrondies relève du \omterm{def:m3:spot-markets:wash}{wash trading}.

\begin{omfigure}
\begin{tikzpicture}
\begin{axis}[width=11cm,height=5cm,ybar,bar width=6pt,
  xlabel={premier chiffre significatif de la taille de transaction},ylabel={part des transactions, \%},
  tick label style={font=\small},label style={font=\small},
  xmin=0.4,xmax=9.6,xtick={1,2,3,4,5,6,7,8,9},ymin=0,ymax=35,grid=major,grid style={gray!25},
  legend columns=3,legend style={font=\footnotesize,at={(0.5,-0.3)},anchor=north,draw=none,fill=none,/tikz/every even column/.append style={column sep=8pt}},
  table/col sep=comma]
\addplot[fill=gray!40,draw=gray,area legend] table[x=digit,y=benford]{figdata/markets-3/16-spot-markets/digits.csv};\addlegendentry{loi de Benford}
\addplot[fill=omDef!60,draw=omDef,area legend] table[x=digit,y=genuine]{figdata/markets-3/16-spot-markets/digits.csv};\addlegendentry{séquence authentique}
\addplot[fill=omThm!60,draw=omThm,area legend] table[x=digit,y=mixed]{figdata/markets-3/16-spot-markets/digits.csv};\addlegendentry{70~\% de wash trades}
\end{axis}
\end{tikzpicture}

\omcaption{Premiers chiffres significatifs de 20\,000 tailles de transactions~: une séquence authentique simulée suit la loi de Benford
(khi-deux de 1.4 face à une valeur critique de 15.5 au seuil de 5~\%)~; une séquence dans laquelle 70~\% des transactions sont produites par une machine, avec des tailles
uniformes entre 0.0010 et 0.0500 bitcoin, ne la suit pas. Séquences synthétiques. Données~: le tutoriel du chapitre.}\label{fig:m3:spot-markets:benford}
\end{omfigure}

\begin{omfigure}
\begin{tikzpicture}
\begin{axis}[width=11cm,height=5cm,
  xlabel={taille de transaction, BTC (intervalles de 0.001)},ylabel={part des transactions, \%},
  tick label style={font=\small},label style={font=\small},
  xmin=0,xmax=0.05,ymin=0,ymax=8,grid=major,grid style={gray!25},
  xticklabel style={/pgf/number format/fixed,/pgf/number format/precision=2},
  scaled x ticks=false,
  legend columns=2,legend style={font=\footnotesize,at={(0.5,-0.3)},anchor=north,draw=none,fill=none,/tikz/every even column/.append style={column sep=8pt}},
  table/col sep=comma]
\addplot[omDef,very thick,const plot] table[x=size,y=genuine]{figdata/markets-3/16-spot-markets/sizes.csv};\addlegendentry{séquence authentique}
\addplot[omThm,very thick,dashed,const plot] table[x=size,y=mixed]{figdata/markets-3/16-spot-markets/sizes.csv};\addlegendentry{70~\% de wash trades}
\end{axis}
\end{tikzpicture}

\omcaption{Tailles de transactions allant jusqu'à 0.05 bitcoin~: les échanges authentiques se regroupent sur des tailles rondes (les intervalles commençant à 0.01,
0.02, 0.03 et 0.04 bitcoin)~; le volume produit par des machines dilue ces regroupements. Séquences synthétiques. Données~: le
tutoriel du chapitre.}\label{fig:m3:spot-markets:sizes}
\end{omfigure}

Le volume publié constitue donc une mesure médiocre de l'endroit où négocier. Ce qui importe à une société qui doit exécuter,
c'est ce qu'elle peut échanger sans faire bouger le prix.

\begin{definition}[Profondeur affichée]\label{def:m3:spot-markets:depth}
\emph{profondeur affichée}\index{profondeur affichée} est la quantité affichée dans un carnet d'ordres à une distance donnée du
prix milieu, par exemple à 10 points de base de chaque côté, sommée sur les niveaux de prix et exprimée dans la
devise de cotation.
\end{definition}

La profondeur peut aussi être truquée, au moyen d'ordres annulés avant que quiconque ne les atteigne~; une société la mesure donc dans la durée,
la confronte à ses propres exécutions et compare les plateformes selon ce qu'il en coûte d'exécuter ses propres tailles~: le glissement réalisé de ses transactions,
la part de ses ordres exécutés au prix affiché et la rapidité avec laquelle le carnet se reconstitue après
une transaction. Ces mesures, et non le volume publié, décident de l'endroit où la société envoie ses ordres et de celui où elle
conserve son inventaire.

\section{Tutoriel~: balayages d'arbitrage et tests de wash trading}\label{tut:m3:spot-markets:scan}

\begin{tutorial}
\textbf{Objectif.} Balayer des carnets synthétiques sur trois plateformes à la recherche d'opportunités triangulaires et entre plateformes après frais, chiffrer
le risque de \omterm{def:m3:spot-markets:latency}{latence de transfert}, et appliquer deux tests de \omterm{def:m3:spot-markets:wash}{wash trading} à des séquences synthétiques. \textbf{État final~:}
les \cref{fig:m3:spot-markets:latency,fig:m3:spot-markets:benford,fig:m3:spot-markets:sizes} et les décomptes du
texte.
\begin{tutsteps}
\item \textbf{Le triangle.} Les deux sens d'un cycle, avec la taille exécutable en haut des carnets.
\omcode{code/firm/triarb/firm_triarb.py}{30}{43}{Les deux sens d'un cycle triangulaire sur une seule plateforme.}\label{lst:m3:spot-markets:tri}
\item \textbf{Le balayage.} \texttt{count\_opportunities(fee)} lance le détecteur sur 500 instantanés, avec des frais preneur de 0,
2, 5 et 10 points de base.
\item \textbf{Les tests.} Premiers chiffres comparés à la loi de Benford, regroupement sur les multiples de 0.01 bitcoin, et
l'estimation de la part de \omterm{def:m3:spot-markets:wash}{wash trading} fondée sur l'arrondi.
\omcode{code/markets-3/16-spot-markets/python/m3_spot.py}{93}{113}{Rapport de regroupement et estimation de la part de wash trades fondée sur l'arrondi.}\label{lst:m3:spot-markets:wash}
\item \textbf{Exécuter} \texttt{fig\_spot.py} pour obtenir les données des graphiques.
\end{tutsteps}
\textbf{Pour aller plus loin.} Faire arrondir ses tailles au trader de \omterm{def:m3:spot-markets:wash}{wash trading} et voir quel test le détecte encore~; ajouter un test de queues épaisses sur
les transactions les plus importantes~; faire suivre aux prix des plateformes synthétiques un même meneur avec un retard et mesurer quelle plateforme
mène.
\end{tutorial}

\section{Construction~: le détecteur d'arbitrages}\label{bld:m3:spot-markets:triarb}

\begin{build}
\bfield{Rôle} Repérer, parmi les plateformes auxquelles la petite société est connectée, les cycles triangulaires et les transactions entre plateformes
qui restent rentables après tous les coûts, et les classer selon le profit espéré~; alimenter la couche d'exécution et le planificateur
d'inventaire.
\bfield{Interface} \texttt{Quote(bid, ask, bid\_size, ask\_size)}~; \texttt{triangular(book, a, b, q, fee)} renvoie
les deux sens~; \texttt{cross(buy, sell, fee\_buy, fee\_sell, route, rebalance\_bp, transfer\_minutes, vol\_annual,
z)}~; \texttt{scan(books, fees, cycles, rebalance\_bp, min\_edge\_bp)} renvoie les opportunités, la plus rentable en premier.
\bfield{Règles} Frais preneur par plateforme (selon le palier de chaque plateforme)~; tailles limitées par le meilleur niveau de chaque carnet sur la route~;
les transactions issues de l'inventaire supportent le coût de rééquilibrage, et les transactions avec transfert supportent $z$ écarts types du
mouvement pendant la durée du transfert (à partir des durées de confirmation de \texttt{firm.gasfee} augmentées du temps de traitement des plateformes).
\bfield{Tests d'acceptation} \texttt{code/firm/triarb/tests/}~: les deux sens triangulaires comparés à un calcul à la main~; des
frais qui suppriment une opportunité~; la pénalité de risque de transfert~; le classement et le filtrage du balayage.
\bfield{Pour aller plus loin} Profondeur au-delà du meilleur niveau du carnet~; cours croisés des \omterm{def:m3:blockchains-for-traders:stable}{stablecoins} valorisés
au prix de marché plutôt qu'au pair~; limites d'inventaire par plateforme et par actif~; balayage sur carnets en direct, avec un taux de réussite mesuré par rapport aux exécutions.
\end{build}

\begin{omsources}
\item I.~Makarov et A.~Schoar, «~Trading and arbitrage in cryptocurrency markets~», \emph{Journal of Financial
Economics} 135(2), 2020, 293--319 (version acceptée, LSE Research Online).
\item L.~W.~Cong, X.~Li, K.~Tang et Y.~Yang, «~Crypto Wash Trading~», NBER Working Paper 30783, 2022~; publié dans
\emph{Management Science}, 2023.
\item SEC, compte rendu d'une réunion avec Bitwise Asset Management, NYSE Arca et Vedder Price, File No.\
SR-NYSEArca-2019-01, 20 mars 2019, avec la présentation de Bitwise.
\item Binance, point de terminaison des informations sur les échanges de l'API spot, 24 septembre 2026.
\item Financial Services Commission (Corée), communiqué de presse sur le trading de cryptomonnaies sous identité réelle, 23 janvier 2018.
\end{omsources}

\section{Exercices}

\begin{exercise}[$\star$]\label{exo:m3:spot-markets:1}
Le BTC/USDT vaut 60\,000 et l'ETH/USDT 3\,000. Quel est le cours croisé ETH/BTC implicite~? À des frais preneur de 10 points de base,
à quelle distance environ le cours ETH/BTC coté peut-il s'écarter de celui-ci avant qu'un \omterm{def:m3:spot-markets:tri}{arbitrage triangulaire} ne rapporte~?
\end{exercise}

\begin{exercise}[$\star$]\label{exo:m3:spot-markets:2}
Le BTC/USDT vaut 60\,030 sur une plateforme et le BTC/USD 60\,000 sur une autre. Quel taux USDT/USD rend les deux prix
cohérents~?
\end{exercise}

\begin{exercise}[$\star$]\label{exo:m3:spot-markets:3}
Pourquoi un écart de prix du bitcoin entre deux pays peut-il être bien plus grand que l'écart du cours ether--bitcoin entre
ces mêmes pays~?
\end{exercise}

\begin{exercise}[$\star\star$]\label{exo:m3:spot-markets:4}
Un écart de 100 points de base, des coûts de 20, une volatilité de 60~\% par an. Quelle est la probabilité de perte si le transfert
dure 30 minutes~? 60 minutes~?
\end{exercise}

\begin{exercise}[$\star\star$]\label{exo:m3:spot-markets:5}
Sur une bourse de référence, 40~\% des transactions d'au moins 0.01~bitcoin portent sur des multiples de 0.01. Sur une autre, 12~\% en portent.
Estimer sa part de \omterm{def:m3:spot-markets:wash}{wash trading} parmi ces transactions.
\end{exercise}

\begin{exercise}[$\star\star$]\label{exo:m3:spot-markets:6}
Expliquer pourquoi un \omterm{def:m3:spot-markets:cross}{inventaire prépositionné} supprime le risque de transfert d'un arbitrage, mais non la totalité de ses risques.
\end{exercise}

\begin{exercise}[$\star\star\star$]\label{exo:m3:spot-markets:7}
\emph{Programmation.} À l'aide de \texttt{count\_opportunities}, déterminer combien d'opportunités triangulaires et entre plateformes rapportent sur 500
instantanés, à des frais preneur de 0, 2, 5 et 10 points de base. Pourquoi les opportunités entre plateformes sont-elles plus nombreuses que les triangles~?
\end{exercise}

\begin{exercise}[$\star\star\star$]\label{exo:m3:spot-markets:8}
\emph{Repérer l'erreur.} «~Cette bourse affiche deux fois le volume de celle-là, donc nos ordres y coûteront moins cher.~»
\end{exercise}

\section{Problème~: la prime kimchi}

\begin{problem}[{Problème du week-end --- tarification d'un contrôle des capitaux}]\label{pb:m3:spot-markets:1}
Un résident de Corée disposant de comptes à l'étranger envisage l'aller-retour suivant~: acheter du bitcoin avec des dollars à l'étranger (frais
preneur de 0.10~\%), le retirer (0.05~\% de la valeur), le transférer en une heure, le vendre contre des wons en Corée (frais preneur de 0.05~\%), convertir
les wons en dollars (fourchette de 0.5~\%) et rapatrier ces dollars à l'étranger par un canal qui coûte 2~\% tout compris. La volatilité est de 70~\%
par an~; le trader facture au transfert 1.65 écart-type du mouvement de prix. Tous les coûts sont fournis à titre illustratif.

\medskip\noindent\textbf{Partie I --- Les coûts.}
\begin{enumerate}
\item Combien coûte un dollar de bitcoin, livré au point de retrait, à l'étranger~?
\item Quelle est la charge de risque de prix du transfert~?
\item Quelle fraction des wons encaissés revient sous forme de dollars à l'étranger~?
\item Quelle prime rend l'aller-retour à l'équilibre~?
\item Avec la prime moyenne de 15~\% observée de décembre 2017 au début de février 2018, quel est le profit espéré par dollar,
hors charge de risque~?
\end{enumerate}

\medskip\noindent\textbf{Partie II --- La contrainte.}
\begin{enumerate}[resume]
\item Quelle serait la prime d'équilibre si les wons pouvaient être rapatriés sans autre coût que la fourchette~?
\item Et si le transfert durait dix minutes au lieu d'une heure~?
\item Pour USD~50\,000, quel est le profit espéré d'un aller-retour à 15~\%~?
\item Pourquoi une société étrangère ne peut-elle pas simplement vendre du bitcoin sur une bourse coréenne~?
\item Pourquoi personne ne pouvait-il vendre à découvert en Corée contre la prime~?
\end{enumerate}

\medskip\noindent\textbf{Partie III --- Les éléments empiriques.}
\begin{enumerate}[resume]
\item Que s'est-il passé pour la prime lorsque le bitcoin a fortement monté, et pourquoi~?
\item Pourquoi l'écart du cours ether--bitcoin entre la Corée et les États-Unis n'était-il que d'environ 3~\%~?
\item Que nous apprennent les 1.275~milliards de dollars de profits d'arbitrage estimés entre la Corée et les États-Unis~?
\item Un \omterm{def:m3:spot-markets:cross}{inventaire prépositionné} en Corée lèverait-il la contrainte~?
\item Transférer la valeur sous forme de \omterm{def:m3:blockchains-for-traders:stable}{stablecoin} plutôt que de wons permettrait-il d'échapper à la contrainte~?
\end{enumerate}

\medskip\noindent\textbf{Partie IV --- Jugement.}
\begin{enumerate}[resume]
\item La \omterm{def:m3:spot-markets:kimchi}{prime kimchi} est-elle une inefficience~?
\item Qui en a profité, et quels risques ont-ils pris~?
\item Que nous apprend la prime sur l'endroit où se forment les prix~?
\item Énoncer le \emph{résultat nommé}~: la prime d'équilibre de l'aller-retour avec et sans coût de
rapatriement.
\item En une phrase~: qu'est-ce qui limite l'\omterm{def:m3:spot-markets:cross}{arbitrage entre plateformes} de crypto-actifs~?
\end{enumerate}
\end{problem}

\section{Questions d'entretien}

\begin{interviewq}[$\star$ \iqroles{trader}]\label{iq:m3:spot-markets:1}
Qu'est-ce que l'\omterm{def:m3:spot-markets:tri}{arbitrage triangulaire}, et pourquoi est-il rare sur une grande plateforme~?
\end{interviewq}

\begin{interviewq}[$\star$ \iqroles{trader, risque}]\label{iq:m3:spot-markets:2}
Pourquoi un cours BTC/USDT n'est-il pas un prix en dollars, et qu'est-ce que cela change pour vos livres~?
\end{interviewq}

\begin{interviewq}[$\star\star$ \iqroles{chercheur}]\label{iq:m3:spot-markets:3}
Comment testeriez-vous si le volume déclaré par une bourse est réel~?
\end{interviewq}

\begin{interviewq}[$\star\star$ \iqroles{trader}]\label{iq:m3:spot-markets:4}
Deux plateformes diffèrent de 50 points de base pendant une heure. Pourquoi personne ne pourrait-il arbitrer cet écart~?
\end{interviewq}

\begin{interviewq}[$\star\star$ \iqroles{risque}]\label{iq:m3:spot-markets:5}
Quels risques porte l'\omterm{def:m3:spot-markets:cross}{inventaire prépositionné} d'un desk d'arbitrage, et comment les limiteriez-vous~?
\end{interviewq}

\begin{interviewq}[$\star\star\star$ \iqroles{développeur, chercheur}]\label{iq:m3:spot-markets:6}
Concevez un système qui décide où conserver l'inventaire sur dix plateformes pour une stratégie d'\omterm{def:m3:spot-markets:cross}{arbitrage entre plateformes}.
\end{interviewq}
